% -------------------------------------------------------------------------
% WBS ID: RES-ML-DAT
% Deliverable: Learning Dataset and Simulation-Data Requirements
% Status: In progress
% Evidence status: ML dataset deferred by current hydrodynamic scope
% Review status: Not reviewed
% Last updated: 2026-07-19
% Dependencies: RES-DAT-OBS, RES-DAT-SCN, RES-DAT-CAS, RES-GEO-SPEC, RES-VER-SPEC
% Downstream interfaces: RES-ML-REP, RES-ML-TRN, RES-ML-SUR
% -------------------------------------------------------------------------

\section{RES-ML-DAT --- Learning Dataset and Simulation-Data Requirements}
\label{sec:res-ml-dat}

\subsection{Purpose and Scope}
\label{sec:res-ml-dat-purpose}

This Deliverable records future learning-dataset requirements only.
The active project builds hydrodynamic source, corridor, geometry and
output manifests; it does not create an ML training dataset. A future
dataset sample would be a complete simulation case, not an extracted
time step. The deferred dataset relationship is:
\begin{align}
  &\text{corridor identity}
  + \text{fixed-rigid barrier geometry}
  + \text{validated tsunami condition}
  + \text{simulation fidelity and numerical configuration}
  \notag\\
  &\qquad \rightarrow
  \text{hydraulic response}
  + \text{uncertainty and quality state}
  \label{eq:res-ml-dat-sample-definition}
\end{align}

Historical pre-event and current or reconstructed geometries may define
realistic intervention features and contextual or weak labels, but they shall
not be treated as direct causal ground truth for optimal reinforcement.

\subsection{Research Questions}
\label{sec:res-ml-dat-research-questions}

% TODO: State the questions that must be answered before the Deliverable
% can be accepted.

\subsection{Terminology and Definitions}
\label{sec:res-ml-dat-definitions}

% TODO: Define the technical terminology, symbols, quantities and
% classifications used in this Deliverable.

\subsection{Literature Evidence}
\label{sec:res-ml-dat-literature-evidence}

% TODO: Record evidence from peer-reviewed papers, authoritative datasets,
% standards, technical reports and primary documentation.
%
% Do not create a detached literature-summary list. Explain how each source
% supports, contradicts or limits a project decision.

\subsection{Governing Relations and Quantitative Definitions}
\label{sec:res-ml-dat-governing-relations}

% TODO: Record relevant equations, dimensionless groups, empirical models,
% extraction rules or quantitative definitions.
%
% Use align environments for displayed equations.
% Every displayed equation must have a unique \label{eq:...}.
% Do not place punctuation after displayed equations.

\subsection{Machine-Learning Role}
\label{sec:res-ml-dat-machine-learning-role}

% TODO: Complete this RES-ML Domain-specific subsection.

\subsection{Non-ML Baseline}
\label{sec:res-ml-dat-non-ml-baseline}

% TODO: Complete this RES-ML Domain-specific subsection.

\subsection{Input and Output Representation}
\label{sec:res-ml-dat-input-output-representation}

% TODO: Complete this RES-ML Domain-specific subsection.

\subsection{Dataset Requirements}
\label{sec:res-ml-dat-dataset-requirements}

Future ML data groups may include:
\begin{itemize}
  \item corridor, geometry and case identity;
  \item fixed-rigid barrier geometry;
  \item local tsunami-condition descriptors;
  \item fluid-response outputs;
  \item simulation metadata;
  \item convergence state;
  \item physical-validity state;
  \item uncertainty state;
  \item future structural, damage or adequacy labels only if those
    extensions are activated and validated.
\end{itemize}

No ML training dataset may be accepted before the selected event,
local hydrodynamic model and impact/loading outputs have passed the
V1--V2--V3 validation gates.

\subsection{Model or Architecture Candidates}
\label{sec:res-ml-dat-model-architecture-candidates}

% TODO: Complete this RES-ML Domain-specific subsection.

\subsection{Training and Validation Method}
\label{sec:res-ml-dat-training-validation-method}

The previous mandatory secondary-event holdout requirement remains
removed from the active project. Future ML validation would require
grouped partitioning based on complete simulation cases, held-out
geometry configurations, held-out condition ranges and independent
final test simulations.

Time steps from the same simulation shall never be randomly split between
future training and test datasets.

\subsection{Evaluation Metrics}
\label{sec:res-ml-dat-evaluation-metrics}

% TODO: Complete this RES-ML Domain-specific subsection.

\subsection{Generalisation and Uncertainty}
\label{sec:res-ml-dat-generalisation-uncertainty}

% TODO: Complete this RES-ML Domain-specific subsection.

\subsection{Simulation and Optimisation Integration}
\label{sec:res-ml-dat-simulation-optimisation-integration}

% TODO: Complete this RES-ML Domain-specific subsection.

\subsection{Candidate Approaches}
\label{sec:res-ml-dat-candidate-approaches}

% TODO: Define the credible candidate approaches considered.

\subsection{Critical Comparison}
\label{sec:res-ml-dat-critical-comparison}

% TODO: Compare candidates using physically and project-relevant criteria.
% Do not select an approach solely on popularity or implementation ease.

\subsection{Assumptions and Applicability}
\label{sec:res-ml-dat-assumptions}

% TODO: State assumptions and the conditions under which they are valid.

\subsection{Limitations and Failure Conditions}
\label{sec:res-ml-dat-limitations}

% TODO: State where the evidence, model or selected approach ceases to be
% reliable or applicable.

\subsection{Project-Specific Conclusions}
\label{sec:res-ml-dat-conclusions}

The initial learning dataset is a controlled simulation-derived dataset for one
validated event-conditioned domain. It is not a broad multi-event dataset and
does not support cross-event generalisation claims.

\subsection{Selected Approach and Rationale}
\label{sec:res-ml-dat-selected-approach}

% TODO: Record the selected approach only after sufficient evidence exists.
% State the alternatives rejected and the reasons for rejection.

\subsection{Unresolved Decisions}
\label{sec:res-ml-dat-unresolved-decisions}

% TODO: Record unresolved questions, required evidence and decision owners.

\subsection{Requirements and Handoffs}
\label{sec:res-ml-dat-requirements}

% TODO: State requirements passed to other Research Domains, Software,
% verification activities or later execution work.

\subsection{Deliverable Completion Record}
\label{sec:res-ml-dat-completion-record}

% TODO: Record whether the Deliverable is:
%
% - Not started
% - In progress
% - Draft complete
% - Reviewed
% - Accepted
%
% Acceptance requires:
%
% - the research questions to be answered;
% - claims to be supported by appropriate evidence;
% - assumptions and limitations to be explicit;
% - the project-specific conclusion to be stated;
% - downstream requirements to be traceable;
% - unresolved decisions to be closed or formally deferred.
